\subsection{Два определения предела функции в точке. Доказательство эквивалентности (из определения по Гейне определение по Коши).}

\subsubsection*{Предел функции в точке}

Пусть функция $f(x)$ определена в проколотой окрестности точки $a$: $\mathring U_a = U_a \setminus \{a\}$.
\[ a \in (d, c) \setminus \{a\} \]

\begin{tcolorbox}[title=Определение 1 (по Коши / на языке $\varepsilon - \delta$)]
	Число $A$ называется пределом функции $f(x)$ при $x \to a$, если:
	\[ \forall \varepsilon > 0\;\; \exists \delta(\varepsilon) > 0: \forall x \in \mathring U_a \quad (0 < |x-a| < \delta) \Rightarrow |f(x) - A| < \varepsilon \]
	\textit{Обозначение:} $\lim_{x \to a} f(x) = A$.
\end{tcolorbox}

\begin{tcolorbox}[title=Определение 2 (по Гейне / на языке последовательностей)]
	Число $A$ называется пределом функции $f(x)$ при $x \to a$, если:
	Для любой последовательности $\{x_n\}$, сходящейся к $a$ ($x_n \neq a$), соответствующая последовательность значений функции $\{f(x_n)\}$ сходится к числу $A$.
	\[ \forall \{x_n\} \subset \mathring U_a: \lim_{n\to\infty} x_n = a \implies \lim_{n\to\infty} f(x_n) = A \]
\end{tcolorbox}

\subsubsection*{Эквивалентность определений}
\begin{tcolorbox}
	\textbf{Теорема.} Определение по Коши $\Leftrightarrow$ Определение по Гейне.
\end{tcolorbox}

\begin{tcolorbox}[title=Доказательство]
	\textbf{2. Гейне $\Rightarrow$ Коши} \\
	Дано: для любой последовательности $x_n \to a$ следует $f(x_n) \to A$.
	
	\textit{Шаг А (Единственность предела):} Докажем, что все последовательности сходятся к \textbf{одному и тому же} числу $A$.
	Предположим противное: $x'_n \to a \Rightarrow f(x'_n) \to A'$, и $x''_n \to a \Rightarrow f(x''_n) \to A''$, где $A' \neq A''$.
	Составим смешанную последовательность:
	\[ \{y_n\} = \{x'_1, x''_1, x'_2, x''_2, \dots \} \]
	Очевидно, $y_n \to a$. Тогда по условию Гейне $f(y_n)$ должна иметь предел. Но подпоследовательность четных членов стремится к $A''$, а нечетных — к $A'$. Значит, предела не существует. \textbf{Противоречие.} Значит $A' = A'' = A$.

	\textit{Шаг Б (Основное доказательство от противного):}
	Пусть предел по Гейне существует, но по Коши — нет.
	Построим отрицание определения Коши:
	\[ \exists \varepsilon > 0: \forall \delta > 0 \;\; \exists x_\delta: 0 < |x_\delta - a| < \delta, \text{ но } |f(x_\delta) - A| \ge \varepsilon \]
	
	Будем брать $\delta = \frac{1}{n}$ для $n=1, 2, \dots$.
	Найдется такая последовательность точек $\{x_n\}$:
	\[ 0 < |x_n - a| < \frac{1}{n}, \quad \text{но} \quad |f(x_n) - A| \ge \varepsilon \]
	
	Из неравенства $|x_n - a| < \frac{1}{n}$ следует, что $x_n \to a$ (по теореме о зажатой посл-ти).
	Тогда по определению Гейне должно быть $f(x_n) \to A$.
	Однако, у нас $|f(x_n) - A| \ge \varepsilon$ для всех $n$, что делает сходимость к $A$ невозможной.
	\textbf{Противоречие.}
\end{tcolorbox}
